% ===================================================================
%  TABELL Nr 13 — Hotellet. — Restaurangen. — Kaféet.
%  Traduction 2026 d'apres texto/io, controlee sur texto/fr et
%  texto/es. Consigne : voir texto/sv/10-tabell-01.tex.
%
%  Deux notes, marquees toutes deux « (*) », et deux appels : celui du
%  bloc 01-2 (la chambre) et celui du bloc 09-1 (le menu).
% ===================================================================

%%K t13-noto-1 noto
\VUnotes{143.2mm}{%
\textit{\VUgras{(*) Se om rummets enskildheter, tabell nr 6.}}%
}

%%K t13-apar-1 apar
\VUsaut{3.0mm}
\VUtitre{15.0pt}{60}{TABELL Nr 13}
\VUsaut{1.6mm}
\VUpk{10.2pt}{40}{Hotellet. --- Restaurangen.}
\VUsaut{2.0mm}
\VUpk{10.2pt}{40}{Kaféet.}
\VUsaut{3.4mm}

%%K t13-01-1 p
1. --- Sedan jag i går kväll hade kommit till Royan, tog jag in på « \textit{Oceanens Hotell}
» som hade rekommenderats mig av en vän.

%%K t13-01-2 p
Man gav mig ett vackert \VUgras{rum} \textsuperscript{(2)}, på andra
\VUgras{våningen} \textsuperscript{(1)} där jag sov mycket gott (*).

%%K t13-02-1 p
2. --- I morse, när jag gick ut ur mitt rum, i vilket \VUgras{tjänsteflickan}
\textsuperscript{(3)} hade burit in frukosten, gick jag in i \VUgras{rökrummet}
\textsuperscript{(5)} som också används som \VUgras{läsrum} \textsuperscript{(4)} för att läsa
\VUgras{tidningarna} \textsuperscript{(6)}, medan jag rökte en \VUgras{cigarrett}
\textsuperscript{(8)}. Efter läsningen gick jag ner med \VUgras{hissen} \textsuperscript{(9)}
och gick och promenerade genom staden.

%%K t13-03-1 p
3. --- Eftersom jag hade glömt att fråga hotellets \VUgras{tjänsteman} \textsuperscript{(12)}
vid vilken timme man serverar måltiderna vid
\VUgras{gemensamma bordet} \textsuperscript{(11)}, kom jag tillbaka tidigt och
begagnade det för att gå och bada i hotellets \VUgras{badrum} \textsuperscript{(10)}. Sedan
gick jag åter till \VUgras{kassan} \textsuperscript{(15)} för att ringa till en av mina
vänner. Men \VUgras{telefonen} \textsuperscript{(13)} var upptagen; när \VUgras{kassörskan}
\textsuperscript{(14)}, som skrev in en \VUgras{räkning} \textsuperscript{(17)} i
\VUgras{resanderegistret} \textsuperscript{(16)}, såg min förlägenhet, bad hon
\VUgras{portieren} \textsuperscript{(18)} att skicka \VUgras{pickolon}
\textsuperscript{(19)} för att uträtta mitt ärende \VUgras{på cykel} \textsuperscript{(20)}.

%%K t13-04-1 p
4. --- Jag tackade henne och gick ut för att ge en drickspenning åt
\VUgras{bäraren} \textsuperscript{(24)} (packbäraren) som från stationen på sin
\VUgras{handkärra} \textsuperscript{(25)} hade fört min \VUgras{hattask}
\textsuperscript{(26)}, \VUgras{min kappsäck} \textsuperscript{(27)} och
\VUgras{min resväska} \textsuperscript{(28)}. \VUgras{Brevbäraren}
\textsuperscript{(21)} som just hade kommit, gav mig ett \VUgras{brev} \textsuperscript{(22)}.

%%K t13-05-1 p
5. --- Jag stannade ännu någon tid på dörrens tröskel, vid \VUgras{brevlådan}
\textsuperscript{(23)}, för att se på trafiken på gatan. \VUgras{Konduktören}
\textsuperscript{(30)} på \VUgras{omnibussen} \textsuperscript{(29)} lastade på sin vagn
\VUgras{kofferten} \textsuperscript{(31)} åt en \VUgras{resande} \textsuperscript{(32)} som
\VUgras{hotellvärden} \textsuperscript{(33)} tog farväl av. Ett ungt \VUgras{telegrambud}
\textsuperscript{(35)} bar utan brådska ett
\VUgras{telegram} \textsuperscript{(34)} till en av hotellets gäster; en
\VUgras{turist} \textsuperscript{(36)} med sin \VUgras{kamera}
\textsuperscript{(38)} över axeln, rådfrågade sin \VUgras{resehandbok} \textsuperscript{(37)}.

%%K t13-tit-1 sub
\VUsaut{4.0mm}
\VUpk{10.2pt}{40}{Frukostdags.}
\VUsaut{3.0mm}

%%K t13-06-1 p
6. --- Framför gallret slumrade ett sorglöst \VUgras{stadsbud} \textsuperscript{(39)} mellan
sin \VUgras{bärkrok} \textsuperscript{(40)} och sin
\VUgras{blanklåda} \textsuperscript{(41)}, medan en ung \VUgras{tvätterska}
\textsuperscript{(42)} som bar en \VUgras{korg} \textsuperscript{(43)} med linne, skrattade åt
den högtidliga minen hos \VUgras{hovmästaren} \textsuperscript{(44)} som stod vid dörren.

%%K t13-07-1 p
7. --- När jag såg \VUgras{kocken} \textsuperscript{(45)}, som kom ut ur sitt
\VUgras{kök} \textsuperscript{(47)} i \VUgras{källarvåningen}
\textsuperscript{(48)} för att skicka en \VUgras{diskpojke} \textsuperscript{(46)} att uträtta
ett ärende, kom jag ihåg att frukosttimmen är nära och jag gick in i restaurangen, medan jag
gick över gården på vilken en
\VUgras{bilförare} \textsuperscript{(50)} ställde in sin \VUgras{bil}
\textsuperscript{(49)}, medan en \VUgras{mekaniker} \textsuperscript{(52)} lagade, enligt
\VUgras{cyklistens} \textsuperscript{(53)} anvisningar, en
\VUgras{fotogentrehjuling} \textsuperscript{(51)} som just hade råkat ut för en
olycka.

%%K t13-08-1 p
8. --- Jag frågade en ung \VUgras{pickolo} \textsuperscript{(54)} som bar
\VUgras{ölglas} \textsuperscript{(55)}, om det gemensamma bordet står under
verandan, och på hans jakande svar tog jag plats vid ett av borden och lät ge mig en utmärkt
måltid.

%%K t13-noto-2 noto
\VUnotes{143.2mm}{%
\textit{\VUgras{(*) Med hjälp av den rättlista som är införd i ordskatten, ska läraren lätt kunna variera nedanstående matsedel.}}%
}

%%K t13-tit-2 sub
\VUsaut{4.0mm}
\VUpk{10.2pt}{40}{En utmärkt frukost.}
\VUsaut{3.0mm}

%%K t13-09-1 p
9. --- Kyparen bar in matsedeln (*) för frukosten och vinlistan till mig. Jag valde och
beställde som \VUgras{förrätt räkor} med \VUgras{smör}, och, som mellanrätt,
\VUgras{komage på Caens vis}. Jag åt ingen \VUgras{fisk}, utan bad om en portion
\VUgras{lammstek}, och, som \VUgras{grönsaksrätt, spenat} med grädde. Sedan, eftersom ingen av
\VUgras{mellanrätterna} behagade mig, lät jag bära in olika efterrätter, \VUgras{ost},
\VUgras{frukt} och \VUgras{skorpor}. Jag tillade dessutom ett utmärkt
Saint-Émilion-\VUgras{vin}, och mycket nöjd med min måltid lämnade jag bordet sedan jag hade
gett \VUgras{kyparen} \textsuperscript{(57)} en vacker \VUgras{drickspenning}
\textsuperscript{(59)}.

%%K t13-10-1 p
10. --- När \VUgras{föreståndaren} \textsuperscript{(60)} såg mig färdig att gå, föreslog han
mig att gå till balkongen för att beundra \VUgras{Oceanens} \textsuperscript{(62)} praktfulla
panorama. I fjärran började man se
\VUgras{fiskebåtar} \textsuperscript{(63)}, \VUgras{segelfartyg}
\textsuperscript{(64)} och en enorm \VUgras{paketbåt} \textsuperscript{(65)} vars svarta
silhuett avtecknar sig mot de vita \VUgras{molnen} \textsuperscript{(67)} vid
\VUgras{horisonten} \textsuperscript{(66)}. En lätt bris fick \VUgras{flaggan}
\textsuperscript{(70)} att fladdra som är fäst ovanför
\VUgras{skylten} \textsuperscript{(69)} på \VUgras{hallen}
\textsuperscript{(68)}.

%%K t13-tit-3 sub
\VUsaut{3.0mm}
\VUpk{10.2pt}{40}{Nöjen.}
\VUsaut{3.0mm}

%%K t13-11-1 p
11. --- Skådespelet var så fängslande, att jag beslöt att i morgon gå upp på kaféets terrass
med en \VUgras{teaterkikare} \textsuperscript{(71)} eller en
\VUgras{kikare} \textsuperscript{(72)}.

%%K t13-12-1 p
12. --- När jag lämnade balkongen, gick jag till hotellets kafé för att svälja en
\VUgras{dryck} \textsuperscript{(73)}, medan jag spelade ett \VUgras{parti biljard}
\textsuperscript{(75)}. Utan att stanna gick jag genom kaférummet i vilket många gäster
spelade \VUgras{schack} \textsuperscript{(78)},
\VUgras{damspel} \textsuperscript{(79)}, \VUgras{domino}
\textsuperscript{(80)}, \VUgras{bräde} \textsuperscript{(81)} eller
\VUgras{kort} \textsuperscript{(82)}, och jag gick direkt upp till
\VUgras{biljardsalen} \textsuperscript{(74)}.

%%K t13-13-1 p
13. --- När partiet var slut, gick jag ner till bottenvåningen och bad kyparen om ett
\VUgras{kuvert} \textsuperscript{(84)}, \VUgras{papper} \textsuperscript{(83)}, ett
\VUgras{frimärke} \textsuperscript{(85)}, en
\VUgras{penna} \textsuperscript{(87)} och \VUgras{bläck} \textsuperscript{(86)}
för att skriva ett brev.

%%K t13-13-2 p
Sedan kallade jag på en \VUgras{droskkusk} \textsuperscript{(92)} vars
\VUgras{vagn} \textsuperscript{(88)} hade stannat framför kaféet. Jag frågade
honom om priset per timme eller per körning, och, när vi var överens om priset, öppnade han
\VUgras{vagnsdörren} \textsuperscript{(89)} för mig och installerade mig. Han steg upp på
\VUgras{kuskbocken} \textsuperscript{(91)} igen, satte snabbt i gång hästarna med en
\VUgras{liten piska} \textsuperscript{(93)} och vi for i väg på en förtjusande tur.
\VUsaut{6.0mm}
\VUfilet{22mm}
